\documentclass[11pt]{article}
\usepackage[a4paper,margin=25mm]{geometry}
\usepackage{amsmath,amssymb,mathtools}
\usepackage{fontspec,unicode-math}
\setmainfont{Latin Modern Roman}
\setsansfont{Latin Modern Sans}
\setmonofont{Latin Modern Mono}
\setmathfont{Latin Modern Math}
\usepackage{xurl,hyperref}
\hypersetup{colorlinks=true,urlcolor=blue,linkcolor=blue}
\newcommand{\sref}[2]{\hyperlink{source:#1:#2}{[S#2]}}
\newcommand{\cataloguescope}[1]{\paragraph{Mathematical question.}#1}
\begin{document}
% UnsolvedMath ID 8700065; source catalogue code AMR-086-0065
\setcounter{section}{8700064}
\section{Question 5.2 — Bugeaud}\label{8700065}
{\small\sffamily UnsolvedMath ID 8700065 · Graph Theory}
\subsection{Definitions and mathematical statement}

\paragraph{Shared notation.}$\mathbb N=\{1,2,\ldots\}$, and $\mathbb Z,\mathbb Q,\mathbb R,\mathbb C$ denote the integers, rationals, reals and complex numbers. Any other conventions are specified locally.


For an algebraic number $\alpha$, choose its primitive irreducible polynomial $f_\alpha\in\mathbb Z[T]$ with positive leading coefficient $a$ and roots $\alpha^{(1)},\ldots,\alpha^{(d)}$. Set
\[
 M(\alpha)=a\prod_{j=1}^d\max\{1,|\alpha^{(j)}|\},\qquad
 \mathrm h(\alpha)=\frac{\log M(\alpha)}d,
\]
and let $H(\alpha)$ be the maximum absolute value of the coefficients of $f_\alpha$. These are respectively Mahler measure, absolute logarithmic height, and usual polynomial height.

Fix $n\ge2$. Define $\mathrm{ZH}_n\subset\mathbb R$ to consist of those $\xi$ for which constants $c_1(\xi),c_2(\xi)>0$ exist satisfying both: every algebraic $\alpha\in\mathbb C$ of degree at most $n$ obeys $|\xi-\alpha|\ge c_2(\xi)H(\alpha)^{-n-1}$, and infinitely many distinct algebraic $\alpha$ of degree at most $n$ obey $|\xi-\alpha|\le c_1(\xi)H(\alpha)^{-n-1}$. The absolute value is complex modulus.
\paragraph{Statement recorded in the supplied source.}
For every $n\ge2$, is
\[
\{\xi\in\mathbb R:[\mathbb Q(\xi):\mathbb Q]=n+1\}\ \subsetneq\ \mathrm{ZH}_n?
\] \sref{8700065}{2}

\subsection{Short English statement}
For each degree bound, is the set having both uniform lower bounds and infinitely many matching upper approximations strictly larger than the real algebraic numbers of the next degree?
\subsection{Sources}
\begin{description}
\item[\hypertarget{source:8700065:1}{S1}] ULAM.AI and the UnsolvedMath Contributors, UnsolvedMath: A Curated Collection of Open Mathematics Problems, record 8700065 (AMR-086-0065). CC BY 4.0. \url{https://huggingface.co/datasets/ulamai/UnsolvedMath}.
\item[\hypertarget{source:8700065:2}{S2}] Michel Waldschmidt, \emph{Open Diophantine Problems}, arXiv:math/0312440v2 (24 January 2004), Question 5.2, p.46; usual-height convention, p.35. \url{https://arxiv.org/abs/math/0312440v2}.
\end{description}
\label{end:8700065}
\end{document}
